\section{SafeSwitch as a State Machine}
\label{sec:s5g_fsm}

\emph{SafeSwitch} is modeled as a distributed finite-state machine, executed locally by the followers \cite{segata2023multi-technology}. The stable states are denoted by \(F_i\), where \(i\) represents the number of interfaces considered available. Each of these states is associated with a specific driving configuration, namely a pair composed of controller and target gap.

Alongside the stable states there are also states \(G_i\), which represent transition phases. The presence of these states is one of the distinctive aspects of \emph{SafeSwitch}: the system recognizes that a change in communication quality must not produce an instantaneous and abrupt change in vehicle behavior. Before consolidating the new operational state, the platoon must be guided toward distance conditions that are compatible with the controller that is about to be used.

\begin{figure}[htbp]
\centering
\includegraphics[width=\textwidth]{fallback-fsm.pdf}
\caption{SafeSwitch mechanism as a state machine}
\label{fig:fallback-fsm}
\end{figure}

When all interfaces are available, the platoon can exploit the highest-performance controller and a shorter gap. In the experimental configuration of the paper, this corresponds to using the PATH controller in the state with three operative interfaces, a Ploeg-type controller in intermediate states, and finally ACC as a fallback when cooperation can no longer be sustained \cite{ploeg2011design,ploeg2015graceful,segata2023multi-technology}.

The logic is bidirectional. In the event of \emph{failure}, the follower moves down toward more conservative states, associated with lower connectivity and larger safety margins. In the event of \emph{recovery}, by contrast, it can move back up toward higher-performance states, provided that the improvement in communication quality is sufficiently stable.

\begin{figure}[htbp]
\centering
\includegraphics[width=\textwidth]{monitor-fsm.pdf}
\caption{PDR monitoring and SafeSwitch failure/recovery mechanism}
\label{fig:monitor-fsm}
\end{figure}

The distributed nature of the state machine is also relevant. Each follower makes decisions locally on the basis of the messages it receives and of the PDR estimates it computes autonomously. The leader remains conceptually separate from this logic, because it does not have to follow another vehicle in the platoon, whereas the followers are the nodes that directly experience the effect of communication quality on their control law.
